\chapter{Le note estranee all'armonia}
\label{chap:noteestranee}

Questo capitolo tratta tutte le note che appartengono al tessuto melodico ma non all'accordo sottostante. Queste note, dette "estranee" o "non armoniche", arricchiscono la trama musicale rendendo le melodie più fluide e cantabili.

\begin{definizione}[Nota estranea all'armonia]
    \label{def:nota_estranea}
    Una nota si definisce estranea (o non armonica) quando non fa parte dell'accordo in quel momento presente nel basso. Esse creano dissonanze temporanee che risolvono su note reali dell'accordo.
\end{definizione}

\section{Nota di passaggio (NdP)}
\label{sec:nota_passaggio}

\begin{definizione}[Nota di passaggio]
    È una nota estranea che occupa per grado congiunto l'intervallo tra due note reali dell'accordo, consecutive o appartenenti ad accordi diversi.
\end{definizione}

\begin{regola}[Uso della NdP]
    La nota di passaggio cade tipicamente sul tempo debole (o sulla parte debole del tempo). Se posta sul tempo forte, assume un carattere più teso simile all'appoggiatura. Si distingue tra NdP \textbf{diatonica} (nella scala) e \textbf{cromatica} (alterazione).
\end{regola}

\begin{esempio}[Nota di passaggio in Do maggiore]
    Melodia Do-Re-Mi sopra un accordo di Do maggiore: Re è la nota di passaggio.
    \musica{
        \new Staff {
            \key do \major \time 4/4
            \relative do'' {
                do4 \passaggio re mi2 |
            }
        }
    }
\end{esempio}

\section{Nota di volta (o cambiata)}
\label{sec:nota_volta}

\begin{definizione}[Nota di volta]
    È una nota che si allontana dalla nota reale per grado congiunto (superiore o inferiore) e vi ritorna immediatamente.
\end{definizione}

\begin{regola}
    La nota di volta cade solitamente sul tempo debole. La nota di ritorno deve essere identica alla nota reale di partenza.
\end{regola}

\section{Anticipazione}
\label{sec:anticipazione}

\begin{definizione}[Anticipazione]
    È una nota che appartiene all'accordo \textit{successivo} ma appare prima del cambio armonico, sul tempo debole.
\end{definizione}

\begin{regola}
    A differenza del ritardo, l'anticipazione non viene preparata né risolve: essa \textit{è} già la risoluzione anticipata.
\end{regola}

\begin{esempio}[Anticipazione nella cadenza V-I]
    La tonica al soprano appare sull'ultima croma del V7.
    \musica{
        \new Staff {
            \key do \major \time 4/4
            \relative do'' {
                \voiceOne
                <si fa'>4. do8 do2 |
            }
        }
    }
\end{esempio}

\section{Appoggiatura}
\label{sec:appoggiatura}

\begin{definizione}[Appoggiatura]
    È una nota estranea posta sul tempo \textbf{forte}, raggiunta per salto e risolta per grado congiunto discendente (più raramente ascendente).
\end{definizione}

\begin{regola}[Appoggiatura vs Ritardo]
    A differenza del ritardo (\ref{chap:ritardi}), l'appoggiatura non è preparata da una nota tenuta. Storicamente era spesso scritta come abbellimento, ma ha il peso metrico di una nota reale.
\end{regola}

\section{Nota di sfuggita (Échappée)}
\label{sec:sfuggita}

\begin{definizione}[Nota di sfuggita]
    È una nota estranea raggiunta per grado congiunto ma abbandonata per salto in direzione opposta. È comune nel soprano in movimenti melodici veloci.
\end{definizione}

\section{Bordone e nota pedale}
\label{sec:pedale}

\begin{definizione}[Nota pedale]
    È una nota tenuta (solitamente al basso o al soprano) mentre l'armonia cambia sopra o sotto di essa.
\end{definizione}

\begin{regola}
    La nota pedale deve essere consonante con l'accordo all'inizio e alla fine del passaggio. Le dissonanze centrali sono ammesse per l'effetto di accumulo di tensione.
\end{regola}

\section{Sintesi delle note estranee}
\label{sec:sintesi_note_estranee}

\begin{center}
\begin{tabular}{|l|l|l|l|l|}
\hline
Nome & Tempo metrico & Raggiunta per... & Abbandonata per... & Preparata? \\
\hline
Passaggio & Debole & Grado congiunto & Grado congiunto & No \\
Volta & Debole & Grado congiunto & Ritorno & No \\
Anticipazione & Debole & Qualsiasi & Nota reale & No (è già risoluz.) \\
Appoggiatura & Forte & Salto & Grado congiunto & No \\
Sfuggita & Debole & Grado congiunto & Salto opposto & No \\
Pedale & Entrambi & Tenuta & Tenuta & Sì \\
\hline
\end{tabular}
\end{center}
